% UFM-Rec — báo cáo theo mẫu IEEE (conference). Biên dịch: pdflatex ufm_rec_ieee.tex (2 lần).
% Cần IEEEtran.cls và gói vntex (font T5 tiếng Việt). Overleaf có sẵn cả hai.
\documentclass[conference,a4paper]{IEEEtran}
\usepackage[T5]{fontenc}
\usepackage[utf8]{inputenc}
\usepackage{amsmath,amssymb}
\usepackage{graphicx}
\usepackage{booktabs}
\usepackage{multirow}
\usepackage{cite}
\usepackage{xcolor}
\usepackage{tikz}
\usetikzlibrary{positioning,arrows.meta,fit,backgrounds,calc,shapes.geometric}
\usepackage{pgfplots}
\pgfplotsset{compat=1.17}
\usepackage[hidelinks]{hyperref}
\pdfmapfile{+vntex.map}

\renewcommand{\figurename}{Hình}
\renewcommand{\tablename}{Bảng}
\renewcommand{\abstractname}{Tóm tắt}
\renewcommand{\IEEEkeywordsname}{Từ khóa}
\renewcommand{\refname}{Tài liệu tham khảo}

\definecolor{cData}{HTML}{E8EEF9}
\definecolor{cModel}{HTML}{E9F5EC}
\definecolor{cLLM}{HTML}{FDF0E1}
\definecolor{cOut}{HTML}{F3E8F7}
\definecolor{cLine}{HTML}{3B4252}
\tikzset{
  box/.style={draw=cLine, rounded corners=2pt, align=center, font=\scriptsize, inner sep=3pt, minimum height=0.8cm, text width=#1},
  box/.default=2.2cm,
  arr/.style={-{Stealth[length=2mm]}, thick, draw=cLine},
  grp/.style={draw=cLine!60, dashed, rounded corners=4pt, inner sep=5pt},
  lbl/.style={font=\scriptsize\bfseries, text=cLine},
}

\begin{document}

\title{UFM-Rec: Hệ khuyến nghị đa phương thức cho sản phẩm khởi đầu lạnh, mở rộng bằng tín hiệu xu hướng và trợ lý RAG chạy LLM cục bộ được chưng cất}

\author{
\IEEEauthorblockN{Nguyễn Thành Phong, Trần An Kỳ, Nguyễn Hoàng Hải}
\IEEEauthorblockA{Khoa Công nghệ Thông tin, Trường Đại học Văn Lang\\
TP. Hồ Chí Minh, Việt Nam}
}

\maketitle

\begin{abstract}
Sản phẩm mới trên sàn thương mại điện tử thường có đủ tên, mô tả và ảnh nhưng gần như chưa có lượt mua, nên các mô hình dựa trên hành vi khó gợi ý chúng. Bài báo trình bày UFM-Rec, một hệ khuyến nghị kết hợp nhánh hành vi (Transformer nhân quả trên lịch sử) và nhánh ngữ nghĩa (CLIP ViT-B/32 đóng băng) qua trọng số hợp nhất theo độ bất định học được, thực nghiệm trên Amazon Reviews 2023 Toys \& Games (767.045 sản phẩm, 11,57 triệu tương tác, chia theo thời gian). Chúng tôi đóng góp bốn điểm. (1) UFM v2 bổ sung vector LSA của TF-IDF vào nhánh nội dung, cổng hợp nhất biết số tương tác và loss xếp hạng riêng cho từng nhánh, nâng cold macro NDCG@10 từ 0,1855 lên 0,2035 (+9,7\%). (2) Một tín hiệu xu hướng không cần huấn luyện, chỉ dùng dữ liệu train, nâng kết quả kết hợp UFM + TF-IDF từ 0,2207 lên 0,2936 trên nửa validation không dùng để chọn tham số (chênh lệch +0,073, CI95 [0,069; 0,077]). (3) Một trợ lý mua sắm RAG tiếng Việt chạy hoàn toàn cục bộ: LLM lập kế hoạch truy vấn (multi-query, HyDE), truy xuất lai TF-IDF + CLIP gộp bằng RRF, xếp hạng lại bằng cross-encoder, cá nhân hóa bằng UFM v2, đa dạng hóa bằng MMR, tự nới điều kiện khi khớp yếu và kiểm tra bám dữ liệu cho câu trả lời; cross-encoder tăng Recall@10 tìm đúng sản phẩm từ 6,1\% lên 12,2\%. (4) Chưng cất LLM 3B thành 0,5B bằng LoRA trên dữ liệu teacher đã lọc tự động: ở bước viết câu trả lời, mô hình 0,5B đạt 98,8\% câu trả lời có $\geq$2 sản phẩm hợp lệ (teacher 75,0\%) và chỉ 8,9\% lý do chứa số liệu không có trong ngữ cảnh (teacher 11,7\%), với khoảng 1/4 bộ nhớ GPU. Chúng tôi nêu rõ các giới hạn của giao thức đánh giá lấy mẫu.
\end{abstract}

\begin{IEEEkeywords}
hệ khuyến nghị, khởi đầu lạnh, đa phương thức, CLIP, retrieval-augmented generation, chưng cất tri thức, LoRA, LLM cục bộ.
\end{IEEEkeywords}

%=====================================================================
\section{Giới thiệu}
Bài toán khởi đầu lạnh phía sản phẩm \cite{schein2002} xuất hiện khi một sản phẩm có ít hoặc không có tương tác: lọc cộng tác \cite{koren2009,rendle2009} và mô hình chuỗi như SASRec \cite{kang2018}, BERT4Rec \cite{sun2019} không có đủ dữ liệu để học biểu diễn định danh cho nó. Trên tập Toys \& Games, tỷ lệ tương tác rơi vào sản phẩm chưa từng xuất hiện trong train tăng từ 23,3\% ở validation lên 45,7\% ở test, nên năng lực xử lý sản phẩm mới quyết định chất lượng hệ thống khi vận hành lâu dài.

UFM-Rec giải quyết bài toán bằng cách để mô hình tự điều chỉnh mức tin cậy vào nhánh hành vi và nhánh ngữ nghĩa theo từng cặp người dùng -- sản phẩm. Phiên bản đầu (v1) chưa vượt được đối chứng TF-IDF trên nhóm 1--5 tương tác. Bài báo này mô tả vòng cải tiến tiếp theo và một lớp ứng dụng mới:
\begin{itemize}
  \item \textbf{UFM v2}: ba thay đổi kiến trúc và huấn luyện có mục tiêu rõ (mục~\ref{sec:ufm}), cải thiện mọi nhóm sản phẩm.
  \item \textbf{Tín hiệu xu hướng}: điểm độ phổ biến gần đây không cần huấn luyện, kết hợp với UFM và TF-IDF (mục~\ref{sec:trend}).
  \item \textbf{Trợ lý RAG tiếng Việt chạy cục bộ}: hỏi đáp tự nhiên trên 767 nghìn sản phẩm, có trích dẫn và kiểm tra bám dữ liệu (mục~\ref{sec:rag}).
  \item \textbf{Chưng cất 3B $\rightarrow$ 0,5B}: dữ liệu teacher được chọn bằng phản hồi truy xuất và bộ lọc tự động (mục~\ref{sec:distill}).
\end{itemize}
Mọi tham số được chọn trên nửa sớm của validation và báo cáo trên nửa muộn; tập test không được mở trong vòng này.

%=====================================================================
\section{Công trình liên quan}
\textbf{Khuyến nghị tuần tự và đa phương thức.} SASRec \cite{kang2018} và BERT4Rec \cite{sun2019} dùng Transformer \cite{vaswani2017} trên chuỗi tương tác. CLIP \cite{radford2021} cung cấp biểu diễn chung cho ảnh và văn bản, thường được dùng làm đặc trưng nội dung cho sản phẩm lạnh. Ước lượng độ bất định \cite{kendall2017} gợi ý cách cân nhắc các nguồn tín hiệu không đồng đều.

\textbf{RAG nâng cao.} RAG \cite{lewis2020} ghép truy xuất với sinh văn bản. HyDE \cite{gao2023hyde} sinh tài liệu giả định để truy xuất; reciprocal rank fusion (RRF) \cite{cormack2009} gộp nhiều danh sách xếp hạng; cross-encoder \cite{nogueira2019} xếp hạng lại chính xác hơn bi-encoder; MMR \cite{carbonell1998} cân bằng liên quan và đa dạng; Corrective RAG \cite{yan2024crag} sửa truy xuất khi độ khớp thấp.

\textbf{Chưng cất và tinh chỉnh hiệu quả.} Chưng cất tri thức \cite{hinton2015} chuyển năng lực từ mô hình lớn sang mô hình nhỏ; LoRA \cite{hu2022lora} chỉ học ma trận hạng thấp. Chúng tôi dùng họ Qwen2.5 \cite{qwen2024} và học từ đầu ra teacher đã được lọc, tương tự hướng tự huấn luyện có chọn lọc.

%=====================================================================
\section{Dữ liệu và giao thức đánh giá}
\textbf{Dữ liệu.} Amazon Reviews 2023 \cite{hou2024} nhóm Toys \& Games. Tương tác dương là lượt mua đã xác minh từ 4 sao; sau khử trùng lặp còn 11.572.689 tương tác của 6.158.534 người dùng trên 767.045 sản phẩm (gộp theo parent ASIN). Dữ liệu được chia theo thời gian toàn cục 80/10/10 (train trước 11/08/2021, test từ 16/07/2022). Mỗi sản phẩm có văn bản (tên, đặc điểm, mô tả, danh mục) và một ảnh chính; CLIP phủ 99,99\% văn bản và 98,91\% ảnh.

\textbf{Nhóm sản phẩm.} Theo số tương tác train: \emph{zero-shot} (0), \emph{extreme cold} (1--5), \emph{cold} (6--20), \emph{warm} ($>$20).

\textbf{Giao thức v3.} 40.000 ca validation, cân bằng 10.000 ca mỗi nhóm; mỗi ca có sản phẩm đích và 99 mẫu âm lấy theo độ phổ biến trong cùng nhóm, loại mọi sản phẩm dương trước đó của người dùng. Chỉ số chính là \emph{cold macro NDCG@10}, trung bình NDCG@10 \cite{jarvelin2002} của ba nhóm lạnh. Để chọn tham số không thiên lệch, validation được chia theo thời gian thành nửa sớm (\emph{fit}, 20.000 ca) và nửa muộn (\emph{audit}, 20.000 ca). Đây là chỉ số lấy mẫu \cite{krichene2020}, không tương đương xếp hạng toàn catalog.

%=====================================================================
\section{Kiến trúc hệ thống}
Hình~\ref{fig:arch} tóm tắt toàn hệ thống gồm ba tầng: (i) tầng dữ liệu và đặc trưng chạy ngoại tuyến; (ii) tầng khuyến nghị (UFM v2, TF-IDF, xu hướng); (iii) tầng trợ lý hội thoại RAG dùng lại các thành phần của tầng (ii).

\begin{figure*}[t]
\centering
\resizebox{\textwidth}{!}{%
\begin{tikzpicture}[x=1cm,y=1cm]
% (i) dữ liệu
\node[box=2.5cm, fill=cData] (raw)   at (0,0)    {Amazon Reviews 2023\\Toys \& Games\\16,3 tr. đánh giá};
\node[box=2.5cm, fill=cData] (clean) at (3.4,0)  {Lọc mua xác minh $\geq$4 sao,\\khử trùng lặp,\\chia theo thời gian};
\node[box=2.5cm, fill=cData] (feat)  at (6.8,0)  {Đặc trưng sản phẩm:\\TF-IDF 20k, LSA 256,\\CLIP văn bản + ảnh};
\node[box=2.5cm, fill=cData] (hist)  at (10.2,0) {Lịch sử 20 sản phẩm,\\đồ thị dương train,\\số tương tác train};
\node[box=2.5cm, fill=cData] (facts) at (13.6,0) {Bảng thông tin RAG:\\giá, đánh giá,\\độ tuổi, danh mục};
% (ii) khuyến nghị
\node[box=2.5cm, fill=cModel] (ufm)   at (0,-2.3)    {\textbf{UFM v2}\\hành vi + ngữ nghĩa,\\hợp nhất theo độ bất định};
\node[box=2.5cm, fill=cModel] (tfidf) at (3.4,-2.3)  {\textbf{TF-IDF}\\cosine hồ sơ -- ứng viên};
\node[box=2.5cm, fill=cModel] (trend) at (6.8,-2.3)  {\textbf{Xu hướng}\\$\log(1{+}n_D)-\beta\log(1{+}n)$};
\node[box=2.5cm, fill=cModel] (blend) at (10.2,-2.3) {\textbf{Kết hợp}\\$\alpha z_{U}{+}(1{-}\alpha)z_{T}{+}\gamma z_{r}$};
\node[box=2.5cm, fill=cOut]   (recs)  at (13.6,-2.3) {Danh sách gợi ý\\top-$K$};
% (iii) RAG — hai hàng
\node[box=2.5cm, fill=cOut]   (user) at (0,-4.7)    {Câu hỏi tiếng Việt\\(+ lịch sử mua, tùy chọn)};
\node[box=2.5cm, fill=cLLM]   (plan) at (3.4,-4.7)  {\textbf{LLM 1,5B lập kế hoạch}\\truy vấn chính + 2 phụ\\+ HyDE; regex giá/tuổi};
\node[box=2.5cm, fill=cModel] (ret)  at (6.8,-4.7)  {\textbf{Truy xuất lai}\\TF-IDF + CLIP mỗi truy vấn\\$\rightarrow$ RRF, lọc cứng, top-100};
\node[box=2.5cm, fill=cModel] (rr)   at (10.2,-4.7) {\textbf{Xếp hạng lại}\\cross-encoder + phổ biến\\+ UFM v2 (cá nhân hóa)};
\node[box=2.5cm, fill=cModel] (mmr)  at (13.6,-4.7) {\textbf{MMR} chọn 6 đa dạng;\\\textbf{CRAG} nới điều kiện\\khi khớp yếu};
\node[box=2.5cm, fill=cModel] (cmp)  at (13.6,-6.6) {Rút gọn ngữ cảnh\\theo truy vấn};
\node[box=2.5cm, fill=cLLM]   (ans)  at (10.2,-6.6) {\textbf{LLM 0,5B chưng cất}\\trả lời JSON:\\lựa chọn + lý do};
\node[box=2.5cm, fill=cModel] (gr)   at (6.8,-6.6)  {Kiểm tra bám dữ liệu:\\điền tên/giá từ dữ liệu,\\thay lý do có số bịa};
\node[box=2.5cm, fill=cOut]   (ui)   at (3.4,-6.6)  {Giao diện (NDJSON):\\thẻ sản phẩm, trích dẫn,\\bảng so sánh, gợi ý tiếp};
\begin{scope}[on background layer]
\node[grp, fit=(raw)(facts), label={[lbl, anchor=south west]north west:(i) Dữ liệu và đặc trưng (ngoại tuyến)}] (g1) {};
\node[grp, fit=(ufm)(recs), label={[lbl, anchor=south west]north west:(ii) Tầng khuyến nghị}] (g2) {};
\node[grp, fit=(user)(mmr)(ui)(cmp), label={[lbl, anchor=south west]north west:(iii) Trợ lý hội thoại RAG (LLM chạy cục bộ)}] (g3) {};
\end{scope}
\foreach \u/\v in {raw/clean,clean/feat,feat/hist,hist/facts,ufm/tfidf,tfidf/trend,trend/blend,blend/recs,user/plan,plan/ret,ret/rr,rr/mmr,mmr/cmp,cmp/ans,ans/gr,gr/ui} \draw[arr] (\u) -- (\v);
\draw[arr, dashed, draw=cLine!70] (g1.south -| trend.north) -- (trend.north);
\draw[arr, dashed, draw=cLine!70] (g1.south -| ufm.north) -- (ufm.north);
\draw[arr, dashed, draw=cLine!70] (g1.south -| tfidf.north) -- (tfidf.north);
\draw[arr, dashed, draw=cLine!70] (g2.south -| rr.north) -- (rr.north);
\draw[arr, dashed, draw=cLine!70] (g2.south -| ret.north) -- (ret.north);
\end{tikzpicture}}
\caption{Kiến trúc tổng thể UFM-Rec. Tầng (i) dựng dữ liệu và đặc trưng một lần; tầng (ii) chấm điểm ứng viên cho người dùng; tầng (iii) dùng lại bảng TF-IDF/CLIP (mũi tên đứt vào khối truy xuất) và UFM v2 (vào khối xếp hạng lại) để trả lời câu hỏi tự nhiên. Ô cam: bước dùng LLM; ô xanh lá: thành phần không sinh văn bản; ô tím: đầu vào/đầu ra.}
\label{fig:arch}
\end{figure*}

%---------------------------------------------------------------------
\subsection{Mô hình UFM v2}\label{sec:ufm}
UFM cho điểm từng cặp người dùng $u$ -- ứng viên $j$ bằng hai nhánh (Hình~\ref{fig:ufm}).

\textbf{Nhánh ngữ nghĩa.} Với sản phẩm $i$, vector CLIP văn bản và ảnh $\mathbf{e}_i^{t},\mathbf{e}_i^{v}\in\mathbb{R}^{512}$, mặt nạ khả dụng $m_i$ và (mới ở v2) vector LSA $\mathbf{l}_i\in\mathbb{R}^{256}$ của TF-IDF đi qua adapter:
\begin{equation}
\mathbf{v}_i=\mathrm{LN}\big(W_2\,\mathrm{GELU}(W_1[\mathbf{e}_i^{t}\odot m_i^{t},\mathbf{e}_i^{v}\odot m_i^{v},\mathbf{l}_i,m_i])\big).
\end{equation}
LSA bù cho giới hạn 77 token của CLIP: TF-IDF đọc toàn bộ tên, đặc điểm và mô tả. Hồ sơ ngữ nghĩa $\mathbf{z}_u^{s}$ là trung bình chuẩn hóa của $\mathbf{v}_i$ trong lịch sử; vector cặp $\mathbf{h}^{s}_{uj}=\mathrm{MLP}([\mathbf{z}_u^{s},\mathbf{v}_j,\mathbf{z}_u^{s}\odot\mathbf{v}_j,m_j])$.

\textbf{Nhánh hành vi.} ID của các sản phẩm đã biết trong lịch sử (tối đa 20) qua Transformer nhân quả 2 lớp, 4 đầu, $d=128$; ID dropout 0,5 khi huấn luyện. Vector cặp $\mathbf{h}^{c}_{uj}$ ghép biểu diễn người dùng, embedding ứng viên, tích từng phần tử, $\log(1{+}|S_u|)$ và $\log(1{+}c_j)$ với $c_j$ là số tương tác train.

\textbf{Hợp nhất theo độ bất định có cổng số tương tác.} Mỗi nhánh $k\in\{c,s\}$ dự đoán độ bất định $u^{k}_{uj}=\mathrm{Softplus}(\cdot)$. Ở v2, một MLP nhỏ trên $(\log(1{+}c_j),\log(1{+}|S_u|))$ dịch chuyển độ bất định, $\tilde u^{k}_{uj}=u^{k}_{uj}+g^{k}(c_j,|S_u|)$, để mô hình học cách bớt tin nhánh hành vi ở sản phẩm hiếm. Trọng số và điểm:
\begin{gather}
w^{k}_{uj}=\frac{\mathbb{I}(k\in A_{uj})\,e^{-\tilde u^{k}_{uj}}}{\sum_{l\in A_{uj}}e^{-\tilde u^{l}_{uj}}},\\
s_{uj}=\mathrm{MLP}_{\text{out}}\big(w^{c}_{uj}\mathbf{h}^{c}_{uj}+w^{s}_{uj}\mathbf{h}^{s}_{uj}+\mathbf{g}_{uj}\big),
\end{gather}
với $A_{uj}$ là tập nhánh khả dụng (ứng viên zero-shot chỉ có nhánh ngữ nghĩa) và $\mathbf{g}_{uj}$ là thành phần căn chỉnh chéo có cổng.

\textbf{Hàm mất mát.} Ngoài loss xếp hạng dạng BPR \cite{rendle2009} $\mathcal{L}_{\text{rank}}=\frac{1}{|\mathcal{N}|}\sum_{j\in\mathcal{N}}\mathrm{Softplus}(s_{uj}-s_{up})$ và các thành phần căn chỉnh, độ bất định, hiệu chuẩn của v1, v2 thêm loss xếp hạng trên điểm riêng của từng nhánh:
\begin{equation}
\mathcal{L}=\mathcal{L}_{\text{rank}}+0{,}01\mathcal{L}_{\text{align}}+0{,}01\mathcal{L}_{\text{unc}}+0{,}1\mathcal{L}_{\text{cal}}+0{,}1\!\sum_{k}\mathcal{L}^{k}_{\text{rank}},
\end{equation}
để điểm từng nhánh (và phần dư dùng làm nhãn cho độ bất định) thực sự được giám sát. Mỗi mẫu dương có 8 mẫu âm ngẫu nhiên đều và (mới) 8 mẫu âm khó lấy theo độ phổ biến trong cùng nhóm; AdamW \cite{loshchilov2019}, lr $3\cdot10^{-4}$ với warmup 1.000 bước rồi giảm cosine, batch 256, 5 epoch.

\begin{figure}[t]
\centering
\resizebox{\columnwidth}{!}{%
\begin{tikzpicture}[x=1cm,y=1cm]
\node[box=1.6cm, fill=cData]  (clip) at (0,1.1)   {CLIP văn bản\\+ ảnh (đóng băng)};
\node[box=1.6cm, fill=cData]  (lsa)  at (0,-0.1)  {LSA 256 của\\TF-IDF \textit{(v2)}};
\node[box=1.6cm, fill=cData]  (ids)  at (0,-1.7)  {ID lịch sử\\+ vị trí};
\node[box=1.7cm, fill=cModel] (ad)   at (2.4,0.5) {Adapter $\rightarrow\mathbf{v}_i$,\\hồ sơ $\mathbf{z}^s_u$};
\node[box=1.7cm, fill=cModel] (tr)   at (2.4,-1.7){Transformer\\nhân quả $\rightarrow\mathbf{z}^c_u$};
\node[box=1.5cm, fill=cModel] (hs)   at (4.7,0.5) {$\mathbf{h}^s_{uj}$, $u^s_{uj}$};
\node[box=1.5cm, fill=cModel] (hc)   at (4.7,-1.7){$\mathbf{h}^c_{uj}$, $u^c_{uj}$};
\node[box=1.5cm, fill=cLLM]   (gate) at (4.7,-0.6){Cổng số t.tác \textit{(v2)}\\$c_j,\ |S_u|$};
\node[box=1.7cm, fill=cOut]   (fuse) at (7.0,-0.6){Softmax$(-\tilde u)$\\+ căn chỉnh chéo\\$\rightarrow s_{uj}$};
\draw[arr] (clip) -- (ad); \draw[arr] (lsa) -- (ad); \draw[arr] (ids) -- (tr);
\draw[arr] (ad) -- (hs); \draw[arr] (tr) -- (hc);
\draw[arr] (hs.east) -- (fuse.north west); \draw[arr] (hc.east) -- (fuse.south west); \draw[arr] (gate) -- (fuse);
\end{tikzpicture}}
\caption{Mô hình UFM v2. Hai thay đổi kiến trúc so với v1 được in nghiêng: vector LSA trong nhánh ngữ nghĩa và cổng dịch chuyển độ bất định theo số tương tác.}
\label{fig:ufm}
\end{figure}

%---------------------------------------------------------------------
\subsection{Tín hiệu xu hướng và kết hợp điểm}\label{sec:trend}
Gọi $n_D(j)$ là số lượt mua của $j$ trong $D$ ngày cuối của train và $n(j)$ là tổng lượt mua train. Điểm xu hướng
\begin{equation}
r_j=\log(1+n_D(j))-\beta\log(1+n(j))
\end{equation}
ưu tiên sản phẩm đang tăng tốc thay vì sản phẩm phổ biến lâu năm. Với người dùng chưa có lịch sử, hệ thống chỉ dùng $r_j$ ($D=90$, $\beta=0{,}75$). Với người dùng có lịch sử, điểm cuối là $\alpha z(s^{\text{UFM}})+(1-\alpha)z(s^{\text{TF}})+\gamma z(r)$ với $z(\cdot)$ chuẩn hóa trong từng ca; trên nửa fit chọn được $D=180$, $\beta=1{,}5$, $\alpha=0{,}8$, $\gamma=0{,}8$. Tín hiệu chỉ dùng dữ liệu train nên không rò rỉ nhãn.

%---------------------------------------------------------------------
\subsection{Trợ lý RAG chạy LLM cục bộ}\label{sec:rag}
Người dùng mô tả nhu cầu bằng tiếng Việt (người nhận, tuổi, sở thích, ngân sách); hệ thống trả lời theo luồng sự kiện NDJSON để sản phẩm hiện trước, câu trả lời đến sau.

\textbf{Hiểu câu hỏi.} Một lời gọi LLM (Qwen2.5-Instruct \cite{qwen2024}) sinh JSON gồm truy vấn tiếng Anh chính, hai truy vấn phụ, một mô tả sản phẩm giả định (HyDE \cite{gao2023hyde}), từ khóa loại trừ và cờ ``tinh chỉnh câu trước'' cho hội thoại nhiều lượt. Giá (đổi VND sang USD), tuổi và mức đánh giá tối thiểu được tách bằng biểu thức chính quy để không phụ thuộc vào LLM.

\textbf{Truy xuất lai.} Mỗi văn bản truy vấn được chấm trên toàn bộ 767.045 sản phẩm bằng tổ hợp TF-IDF và CLIP văn bản/ảnh (truy vấn văn bản so với vector ảnh); các danh sách được gộp bằng RRF \cite{cormack2009} với $k=60$, sau đó áp bộ lọc cứng (giá, tuổi, đánh giá, loại trừ) và lấy 100 ứng viên.

\textbf{Xếp hạng lại và đa dạng.} Điểm cuối $=z(\text{RRF})+4\,z(\text{CE})+0{,}5\,z(\text{prior})+0{,}5\,z(\text{UFM})$, trong đó CE là cross-encoder ms-marco-MiniLM-L6 \cite{nogueira2019,wang2020minilm}, prior là độ phổ biến 180 ngày gần nhất và UFM là điểm UFM v2 khi người dùng cung cấp lịch sử mua. MMR \cite{carbonell1998} với $\lambda=0{,}8$ chọn 6 sản phẩm. Nếu điểm CE cao nhất dưới $-2$ (khớp yếu), hệ thống lần lượt nới mức đánh giá, giá ($\times1{,}5$) rồi tuổi, và báo rõ cho người dùng (corrective RAG \cite{yan2024crag}).

\textbf{Sinh câu trả lời có kiểm soát.} Văn bản sản phẩm được rút gọn theo các đoạn trùng từ với truy vấn. LLM chỉ trả về JSON gồm câu mở đầu, danh sách lựa chọn (số thứ tự và lý do) và mẹo; \emph{tên và giá được điền từ dữ liệu}, không do LLM viết. Lý do nào chứa con số không có trong ngữ cảnh của sản phẩm đó bị thay bằng mô tả an toàn; trích dẫn ngoài phạm vi bị loại. Giao diện hiển thị thẻ sản phẩm, trích dẫn bấm được, bảng so sánh và gợi ý hỏi tiếp.

%---------------------------------------------------------------------
\subsection{Chưng cất LLM 3B thành 0,5B}\label{sec:distill}
\textbf{Sinh dữ liệu.} Teacher Qwen2.5-3B sinh 445 câu hỏi tiếng Việt từ danh mục và đặc điểm (không lộ tên) của các sản phẩm \emph{không trùng} với sản phẩm dùng để đánh giá. Với bước lập kế hoạch, teacher sinh 2 phương án cho mỗi câu; chúng tôi chạy truy xuất thật cho từng phương án và giữ phương án truy xuất trúng danh mục và sản phẩm gốc tốt hơn (\emph{lọc bằng phản hồi truy xuất}), được 554 mẫu. Với bước trả lời, 380 mẫu được sinh và chỉ giữ 249 mẫu qua bộ lọc: JSON hợp lệ, $\geq$2 lựa chọn hợp lệ, lý do bằng tiếng Việt và mọi con số có trong ngữ cảnh.

\textbf{Huấn luyện.} Học sinh Qwen2.5-0.5B-Instruct với LoRA \cite{hu2022lora} hạng 16, $\alpha=32$ trên mọi phép chiếu attention và MLP (8,8 triệu tham số học được), loss chỉ tính trên JSON đầu ra, một adapter cho cả hai bước. AdamW lr $2\cdot10^{-4}$, batch hiệu dụng 16, 3 epoch, gradient checkpointing; giữ 40 mẫu để chọn epoch. Adapter (35 MB) được gộp vào trọng số khi phục vụ.

%=====================================================================
\section{Thực nghiệm}
\subsection{Môi trường}
Một GPU 24 GB dùng chung với các tác vụ khác (thường chỉ còn 3--7 GB trống), CPU 40 luồng. Mọi mô hình HF được ghim phiên bản; không dùng mã từ xa.

\subsection{UFM v2}
Bảng~\ref{tab:ufm} so sánh NDCG@10 trên toàn bộ validation v3 (checkpoint tốt nhất mỗi lượt). UFM v2 tăng cold macro 9,7\% và cải thiện cả bốn nhóm, rõ nhất ở zero-shot (+0,026) và warm (+0,061).

\begin{table}[t]
\caption{NDCG@10 của UFM v1 và v2 trên validation v3 (40.000 ca)}
\label{tab:ufm}
\centering
\begin{tabular}{lccccc}
\toprule
 & Cold macro & Zero-shot & 1--5 & 6--20 & Warm\\
\midrule
UFM v1 (ep.~4) & 0,1855 & 0,158 & 0,201 & 0,198 & 0,101\\
UFM v2 (ep.~3) & \textbf{0,2035} & \textbf{0,184} & \textbf{0,215} & \textbf{0,212} & \textbf{0,162}\\
\bottomrule
\end{tabular}
\end{table}

\subsection{Kết hợp điểm và tín hiệu xu hướng}
Bảng~\ref{tab:trend} và Hình~\ref{fig:regime} báo cáo trên nửa audit. Kết hợp UFM v2 + TF-IDF ($\alpha=0{,}7$) vượt từng thành phần; thêm tín hiệu xu hướng tăng thêm 0,073 (CI95 bootstrap [0,069; 0,077]) và 0,090 so với TF-IDF. Mức tăng tập trung ở nhóm extreme cold, cold và warm; nhóm zero-shot gần như không đổi vì sản phẩm chưa có lượt mua train thì cũng không có xu hướng. Với người dùng chưa có lịch sử, xu hướng nâng NDCG@10 từ 0,043 (độ phổ biến train) lên 0,147.

\begin{table}[t]
\caption{Cold macro NDCG@10 trên nửa audit (20.000 ca); tham số chọn trên nửa fit}
\label{tab:trend}
\centering
\begin{tabular}{lcc}
\toprule
Phương pháp & UFM v1 & UFM v2\\
\midrule
TF-IDF & 0,2037 & 0,2037\\
UFM & 0,1930 & 0,2115\\
UFM + TF-IDF & 0,2158 & 0,2207\\
UFM + TF-IDF + xu hướng & 0,2876 & \textbf{0,2936}\\
\bottomrule
\end{tabular}
\end{table}

\begin{figure}[t]
\centering
\begin{tikzpicture}
\begin{axis}[
  ybar, bar width=4.2pt, width=\columnwidth, height=4.6cm,
  ymin=0, ymax=0.4, ylabel={NDCG@10}, ylabel style={font=\scriptsize},
  symbolic x coords={Zero-shot,Extreme cold,Cold,Warm}, xtick=data,
  tick label style={font=\scriptsize}, yticklabel style={/pgf/number format/fixed,/pgf/number format/precision=1},
  legend style={font=\scriptsize, at={(0.5,1.02)}, anchor=south, legend columns=4, draw=none, fill=none, /tikz/every even column/.append style={column sep=4pt}},
  enlarge x limits=0.15]
\addplot[fill=gray!40] coordinates {(Zero-shot,0.1869) (Extreme cold,0.2254) (Cold,0.1989) (Warm,0.1463)};
\addplot[fill=green!45!black!35] coordinates {(Zero-shot,0.1932) (Extreme cold,0.2276) (Cold,0.2138) (Warm,0.1729)};
\addplot[fill=blue!35] coordinates {(Zero-shot,0.1987) (Extreme cold,0.2437) (Cold,0.2196) (Warm,0.1717)};
\addplot[fill=orange!70] coordinates {(Zero-shot,0.1994) (Extreme cold,0.3575) (Cold,0.3238) (Warm,0.2748)};
\legend{TF-IDF, UFM v2, UFM v2 + TF-IDF, + xu hướng}
\end{axis}
\end{tikzpicture}
\caption{NDCG@10 theo nhóm sản phẩm trên nửa audit.}
\label{fig:regime}
\end{figure}

\subsection{Truy xuất của trợ lý RAG}
\textbf{Bộ câu hỏi tổng hợp.} 296 câu hỏi tiếng Việt do LLM sinh từ danh mục và đặc điểm của một sản phẩm đích (không lộ tên), kèm ràng buộc giá/tuổi khi có. Nhiệm vụ là tìm lại đúng sản phẩm đích trong 767 nghìn sản phẩm; chúng tôi đo Recall@10 và tỷ lệ 6 gợi ý cùng danh mục lá với sản phẩm đích. Trọng số được chọn trên 148 câu đầu, báo cáo trên 148 câu sau (Bảng~\ref{tab:ragret}).

\begin{table}[t]
\caption{Truy xuất trên 148 câu hỏi audit}
\label{tab:ragret}
\centering
\begin{tabular}{lcc}
\toprule
Cấu hình & Recall@10 & Đúng danh mục@6\\
\midrule
Chỉ TF-IDF & 0,7\% & 11,7\%\\
TF-IDF + CLIP (trọng số đã chỉnh) & 6,1\% & 26,9\%\\
\quad + cross-encoder MiniLM & \textbf{12,2\%} & \textbf{37,8\%}\\
\midrule
\multicolumn{3}{l}{\emph{Mở rộng truy vấn, cùng LLM 1,5B:}}\\
Chỉ truy vấn chính & -- & 33,7\%\\
Multi-query (trung bình) & -- & 34,2\%\\
Multi-query + HyDE (trung bình) & -- & 34,7\%\\
Multi-query + HyDE + RRF & 11,5\% & 35,3\%\\
\bottomrule
\end{tabular}
\end{table}

CLIP đóng góp phần lớn vì truy vấn mô tả đặc điểm trực quan và công dụng hơn là từ khóa trùng tên; cross-encoder tăng gấp đôi Recall@10 với chi phí 0,09 giây mỗi câu (bge-reranker-base cho chất lượng tương đương nhưng chậm gấp đôi). Ba dòng đầu dùng kế hoạch của LLM 3B; với LLM 1,5B, mỗi kỹ thuật mở rộng đều tăng nhẹ tỷ lệ đúng danh mục, và RRF tốt nhất trên nửa fit nên được chọn.

\subsection{Câu trả lời và chưng cất}
Bảng~\ref{tab:distill} so sánh bốn LLM. Bước lập kế hoạch được đo trên cả 296 câu; bước trả lời trên 80 câu audit với \emph{cùng một ngữ cảnh} sản phẩm cho mọi mô hình. ``Số bịa'' là tỷ lệ lý do chứa con số không có trong ngữ cảnh của sản phẩm được chọn.

\begin{table}[t]
\caption{So sánh LLM cho hai bước của RAG}
\label{tab:distill}
\centering
\setlength{\tabcolsep}{3.2pt}
\begin{tabular}{lcccc}
\toprule
 & 0,5B & \textbf{0,5B CC} & 1,5B & 3B (teacher)\\
\midrule
\multicolumn{5}{l}{\emph{Lập kế hoạch (296 câu)}}\\
JSON hợp lệ & 89,5\% & 95,6\% & 99,7\% & 97,3\%\\
Đúng danh mục@6 & 26,1\% & 30,4\% & 36,6\% & 38,1\%\\
Recall@10 & 11,8\% & 9,8\% & 16,2\% & 14,9\%\\
\midrule
\multicolumn{5}{l}{\emph{Trả lời (80 câu)}}\\
$\geq$2 lựa chọn hợp lệ & 38,8\% & \textbf{98,8\%} & 51,3\% & 75,0\%\\
Lý do tiếng Việt & 58,8\% & \textbf{98,8\%} & 87,5\% & 95,0\%\\
Lý do có số bịa & 25,2\% & \textbf{8,9\%} & 27,5\% & 11,7\%\\
Trích dẫn sai & 1,8\% & 0\% & 0\% & 0\%\\
Thời gian (trung vị) & 5,8 s & 6,6 s & 7,4 s & 189 s$^\ast$\\
Bộ nhớ GPU & 1,3 GB & 1,3 GB & 3,3 GB & 5,4 GB\\
\bottomrule
\multicolumn{5}{l}{\scriptsize CC: chưng cất. $^\ast$GPU dùng chung bận 100\%; khi rảnh 3B mất 12,5 s.}
\end{tabular}
\end{table}

Ở bước trả lời, mô hình 0,5B sau chưng cất vượt cả teacher về định dạng, ngôn ngữ và độ bám dữ liệu, vì nó chỉ học từ những đầu ra đã qua bộ lọc; loss trên tập giữ lại giảm từ 1,21 xuống 0,82. Ở bước lập kế hoạch, chưng cất cải thiện tính hợp lệ JSON và tỷ lệ đúng danh mục (+4,3 điểm) nhưng Recall@10 giảm nhẹ, và mô hình lớn vẫn tốt hơn: việc viết truy vấn tiếng Anh tốt đòi hỏi tri thức thế giới mà 0,5B khó học từ 554 mẫu. Vì vậy hệ thống phục vụ cấu hình kết hợp: \emph{1,5B lập kế hoạch, 0,5B chưng cất viết câu trả lời} (tổng khoảng 4,8 GB), khoảng 20--24 giây mỗi câu khi GPU bận.

%=====================================================================
\section{Thảo luận và hạn chế}
\begin{itemize}
\item \textbf{Giao thức lấy mẫu.} Mẫu âm của validation v3 lấy theo độ phổ biến đến cuối tập validation, nên sản phẩm ``đang lên'' trong validation được lợi; một phần mức tăng của tín hiệu xu hướng có thể đến từ đặc điểm này dù tín hiệu chỉ dùng dữ liệu train. Kết quả 100 ứng viên không tương đương xếp hạng toàn catalog \cite{krichene2020}.
\item \textbf{Độc lập của validation.} Checkpoint được chọn trên validation; nửa audit không phải tập kiểm tra độc lập hoàn toàn. Tập test chưa được mở trong vòng này và nên được dùng một lần cho cấu hình cuối.
\item \textbf{Đánh giá RAG.} Câu hỏi do LLM sinh, không phải người dùng thật; mẫu nhỏ (148 câu truy xuất, 80 câu trả lời). Kiểm tra bám dữ liệu chỉ bắt được con số, không bắt được khẳng định định tính sai.
\item \textbf{Mô hình nhỏ.} 0,5B đôi khi viết sai chính tả trong câu mở đầu; tên và giá không bị ảnh hưởng vì được điền từ dữ liệu.
\end{itemize}

\section{Kết luận}
UFM-Rec cho thấy ba hướng bổ sung cho nhau trên bài toán khởi đầu lạnh: cải tiến có mục tiêu trong mô hình hợp nhất đa phương thức (UFM v2, +9,7\% cold macro), tín hiệu xu hướng rẻ nhưng mạnh (+0,073 trên kết hợp), và một lớp hội thoại RAG chạy cục bộ biến khuyến nghị thành câu trả lời có dẫn chứng. Chưng cất trên dữ liệu teacher đã lọc giúp mô hình 0,5B viết câu trả lời đáng tin cậy hơn teacher 3B với chi phí nhỏ hơn nhiều. Hướng tiếp theo gồm đánh giá trên tập test và với người dùng thật, mở rộng dữ liệu chưng cất cho bước lập kế hoạch và học trọng số xếp hạng lại từ phản hồi.

\section*{Lời cảm ơn}
Nhóm cảm ơn giảng viên Phan Hồ Viết Trường đã hướng dẫn đồ án và FITLAB đã cung cấp tài nguyên tính toán.

\section*{Mã nguồn và tái lập}
Mã nguồn, cấu hình, lịch sử huấn luyện và các tệp kết quả JSON nằm trong kho mã của nhóm (nhánh \texttt{results/rebuild-v2-rag-20261109}); tóm tắt số liệu ở \texttt{reports/SUMMARY\_20261109.md}.

\begin{thebibliography}{99}
\bibitem{schein2002} A.~I. Schein, A.~Popescul, L.~H. Ungar, and D.~M. Pennock, ``Methods and metrics for cold-start recommendations,'' in \emph{Proc. SIGIR}, 2002, pp. 253--260.
\bibitem{koren2009} Y.~Koren, R.~Bell, and C.~Volinsky, ``Matrix factorization techniques for recommender systems,'' \emph{Computer}, vol.~42, no.~8, pp. 30--37, 2009.
\bibitem{rendle2009} S.~Rendle, C.~Freudenthaler, Z.~Gantner, and L.~Schmidt-Thieme, ``BPR: Bayesian personalized ranking from implicit feedback,'' in \emph{Proc. UAI}, 2009, pp. 452--461.
\bibitem{kang2018} W.-C. Kang and J.~McAuley, ``Self-attentive sequential recommendation,'' in \emph{Proc. ICDM}, 2018, pp. 197--206.
\bibitem{sun2019} F.~Sun \emph{et al.}, ``BERT4Rec: Sequential recommendation with bidirectional encoder representations from transformer,'' in \emph{Proc. CIKM}, 2019, pp. 1441--1450.
\bibitem{vaswani2017} A.~Vaswani \emph{et al.}, ``Attention is all you need,'' in \emph{Proc. NeurIPS}, 2017.
\bibitem{radford2021} A.~Radford \emph{et al.}, ``Learning transferable visual models from natural language supervision,'' in \emph{Proc. ICML}, 2021.
\bibitem{kendall2017} A.~Kendall and Y.~Gal, ``What uncertainties do we need in Bayesian deep learning for computer vision?'' in \emph{Proc. NeurIPS}, 2017.
\bibitem{hou2024} Y.~Hou, J.~Li, Z.~He, A.~Yan, X.~Chen, and J.~McAuley, ``Bridging language and items for retrieval and recommendation,'' arXiv:2403.03952, 2024.
\bibitem{jarvelin2002} K.~J\"arvelin and J.~Kek\"al\"ainen, ``Cumulated gain-based evaluation of IR techniques,'' \emph{ACM Trans. Inf. Syst.}, vol.~20, no.~4, pp. 422--446, 2002.
\bibitem{krichene2020} W.~Krichene and S.~Rendle, ``On sampled metrics for item recommendation,'' in \emph{Proc. KDD}, 2020, pp. 1748--1757.
\bibitem{lewis2020} P.~Lewis \emph{et al.}, ``Retrieval-augmented generation for knowledge-intensive NLP tasks,'' in \emph{Proc. NeurIPS}, 2020.
\bibitem{gao2023hyde} L.~Gao, X.~Ma, J.~Lin, and J.~Callan, ``Precise zero-shot dense retrieval without relevance labels,'' in \emph{Proc. ACL}, 2023.
\bibitem{cormack2009} G.~V. Cormack, C.~L.~A. Clarke, and S.~B\"uttcher, ``Reciprocal rank fusion outperforms Condorcet and individual rank learning methods,'' in \emph{Proc. SIGIR}, 2009, pp. 758--759.
\bibitem{nogueira2019} R.~Nogueira and K.~Cho, ``Passage re-ranking with BERT,'' arXiv:1901.04085, 2019.
\bibitem{wang2020minilm} W.~Wang \emph{et al.}, ``MiniLM: Deep self-attention distillation for task-agnostic compression of pre-trained transformers,'' in \emph{Proc. NeurIPS}, 2020.
\bibitem{carbonell1998} J.~Carbonell and J.~Goldstein, ``The use of MMR, diversity-based reranking for reordering documents and producing summaries,'' in \emph{Proc. SIGIR}, 1998, pp. 335--336.
\bibitem{yan2024crag} S.-Q. Yan, J.-C. Gu, Y.~Zhu, and Z.-H. Ling, ``Corrective retrieval augmented generation,'' arXiv:2401.15884, 2024.
\bibitem{hinton2015} G.~Hinton, O.~Vinyals, and J.~Dean, ``Distilling the knowledge in a neural network,'' arXiv:1503.02531, 2015.
\bibitem{hu2022lora} E.~J. Hu \emph{et al.}, ``LoRA: Low-rank adaptation of large language models,'' in \emph{Proc. ICLR}, 2022.
\bibitem{qwen2024} Qwen Team, ``Qwen2.5 technical report,'' arXiv:2412.15115, 2024.
\bibitem{loshchilov2019} I.~Loshchilov and F.~Hutter, ``Decoupled weight decay regularization,'' in \emph{Proc. ICLR}, 2019.
\end{thebibliography}

\end{document}
